\begin{textsample}{POS Dim 3 – es – Score 23.00 – es/es\_inf\_16.txt}  \label{ex:f3_pos_061}
Campo de Experiências : Corpo, \textbf{Gestos} e Movimentos A exploração corporal, dos \textbf{gestos} e movimentos expressos pela \textbf{criança} é uma forma vital dela conhecer a si e o mundo a sua \textbf{volta}.
No processo de relacionar -se com o mundo, ela utiliza recursos como a corporeidade, as diferentes linguagens e a emoção nas interações e \textbf{brincadeiras}, espontâneas ou vivenciadas nas práticas cotidianas planejadas intencionalmente pelos \textbf{adultos}.
Nas DCNEI’s, as práticas pedagógicas de toda escola devem garantir experiências que : I – promovam conhecimento de si e do mundo por meio da ampliação de experiências \textbf{sensoriais}, expressivas e corporais que possibilitem \textbf{movimentação} ampla, expressão da individualidade e respeito pelos ritmos e \textbf{desejos} da \textbf{criança} ; II – favoreçam a imersão das \textbf{crianças} nas diferentes linguagens e o progressivo domínio por elas de vários gêneros e formas de expressão : gestual, verbal, \textbf{plástica}, \textbf{dramática} e musical.
Elas se comunicam e se expressam no entrelaçamento entre corpo, emoção e por meio das diferentes linguagens, como a música, a dança, o teatro, as \textbf{brincadeiras} de faz de conta.
As \textbf{crianças} conhecem e reconhecem as sensações e funções de seu corpo.
Com seus \textbf{gestos} e movimentos identificam suas potencialidades e seus limites, desenvolvendo, ao mesmo tempo, a consciência sobre o que é seguro e o que pode ser um risco à sua integridade física.
E o constante contato com seus pares, materiais e espaços favorecem todo este desenvolvimento.
Na Educação \textbf{Infantil}, o corpo das \textbf{crianças} ganha centralidade, pois ele é o principal alvo dos \textbf{cuidados} físicos e de práticas pedagógicas orientadas para a emancipação, a autonomia e a liberdade.
Para tanto, é imprescindível ao professor pensar não só no espaço que deve transmitir segurança e confiança para as \textbf{crianças}, mas também sobre o desenvolvimento integral da primeira infância, ou seja, como as \textbf{crianças} se desenvolvem e aprendem, principalmente em se tratando dos \textbf{bebês} e das \textbf{crianças} com necessidades educativas especiais.
Ao adotar os eixos norteadores “ \textbf{brincadeiras} e interações ”, a instituição escolar oportuniza aos \textbf{bebês} e às \textbf{crianças} a construção de um amplo repertório de movimentos, \textbf{gestos}, olhares, \textbf{sons} e \textbf{mímicas} com o corpo, a exploração e vivência com seus pares, descobertas de variados modos de interação, ocupação e uso do espaço.
Permite aos \textbf{bebês} e \textbf{crianças} dominarem \textbf{progressivamente} os movimentos com o corpo, tais como sentar com apoio, rastejar, engatinhar, escorregar, caminhar \textbf{apoiando} -se em berços, mesas e cordas, saltar, escalar, equilibrar -se, correr, dar cambalhotas, alongar -se etc.
É necessário também que, no planejamento curricular, o professor garanta situações que envolvam as linguagens musicais e cênicas ( \textbf{brincar}, dançar, \textbf{dramatizar} ), promovendo situações que envolvem a multiplicidade em sua cultura.
O Campo de Experiências “ Corpo, \textbf{Gestos} e Movimentos ” trata dos objetivos de aprendizagem e desenvolvimento, e devem garantir os direitos de aprendizagem de modo a possibilitar à \textbf{criança} : Conviver - com \textbf{crianças} e \textbf{adultos} experimentando marcas da cultura corporal nos \textbf{cuidados} pessoais, na dança, música, teatro, artes circenses, \textbf{escuta} de histórias e \textbf{brincadeiras}. \textbf{Brincar} - utilizando \textbf{criativamente} o repertório da cultura corporal e do movimento.
Participar - de atividades que envolvem práticas corporais, desenvolvendo autonomia para cuidar de si.
Explorar - amplo repertório de movimentos, \textbf{gestos}, olhares, produção de \textbf{sons} e de \textbf{mímicas}, descobrindo modos de ocupação e de uso do espaço com o corpo.
Expressar - corporalmente emoções e representações tanto nas relações cotidianas como nas \textbf{brincadeiras}, dramatizações, danças, músicas, contação de histórias.
Conhecer -se - nas diversas oportunidades de interações e explorações com seu corpo.
Na BNCC este Campo de Experiências estabelece que : > Com o corpo ( por meio dos sentidos, \textbf{gestos}, movimentos impulsivos ou intencionais, coordenados ou espontâneos ), as \textbf{crianças}, desde cedo, exploram o mundo, o espaço e os objetos do seu entorno, estabelecem relações, expressam -se, \textbf{brincam} e produzem conhecimentos sobre si, sobre o outro, sobre o universo social e cultural, tornando -se, \textbf{progressivamente}, conscientes dessa corporeidade.
Por meio das diferentes linguagens, como a música, a dança, o teatro, as \textbf{brincadeiras} de faz de conta, elas se comunicam e se expressam no entrelaçamento entre corpo, emoção e linguagem.
As \textbf{crianças} conhecem e reconhecem as sensações e funções de seu corpo e, com seus \textbf{gestos} e movimentos, identificam suas potencialidades e seus limites, desenvolvendo, ao mesmo tempo, a consciência sobre o que é seguro e o que pode ser um risco à sua integridade física.
Na Educação \textbf{Infantil}, o corpo das \textbf{crianças} ganha centralidade, pois ele é o partícipe privilegiado das práticas pedagógicas de \textbf{cuidado} físico, orientadas para a emancipação e a liberdade, e não para a submissão.
Assim, a instituição escolar precisa promover oportunidades ricas para que as \textbf{crianças} possam, sempre animadas pelo espírito \textbf{lúdico} e na interação com seus pares, explorar e vivenciar um amplo repertório de movimentos, \textbf{gestos}, olhares, \textbf{sons} e \textbf{mímicas} com o corpo, para descobrir variados modos de ocupação e uso do espaço com o corpo ( tais como sentar com apoio, rastejar, engatinhar, escorregar, caminhar \textbf{apoiando} -se em berços, mesas e cordas, saltar, escalar, equilibrar -se, correr, dar cambalhotas, alongar -se etc. ) > ( BRASIL, 2017, p. 36-37 )

% matched lemmas: adulto, apoiar, bebê, brincadeira, brincar, criança, criativamente, cuidado, desejo, dramatizar, dramático, escuta, gesto, infantil, lúdico, movimentação, mímico, plástico, progressivamente, sensorial, som, volta
\end{textsample}
